\section{Real-Robot Logs and Physical-Consequence Prediction}
\label{app:spot-data}

Section~\ref{sec:results:spot-shadow} tests the consequence interface on real
Spot-with-arm telemetry. Unlike the simulator branches in
Appendix~\ref{app:robot-decision-interface}, each logged window contains only
the future that was actually executed. We therefore specify the replay input
contract before the prediction and objective-conditioned ranking results.

The Spot-with-arm evidence uses logged ROS bags collected in a lab environment under human teleoperation.
C-PHAST is evaluated as a shadow model: it observes telemetry and possible-next-behavior descriptors, predicts typed physical consequences, and writes objective-conditioned action cards beside the teleoperated run.
The data are converted into overlapping windows with a two-second history and a five-second future horizon.
Each row supervises the future that was actually executed in the log.
Thus Spot serves as the real-robot telemetry-to-action-card forecast: the log supplies physical supervision for the executed future, while branch execution and matched-state alternative rollouts are tested by the Isaac experiments.
Table~\ref{tab:spot-dataset-arxiv} gives the scale of this logged dataset and the history/future windows used by the reported result.
Training episodes comprise static arm scanning (105 windows), fast traversal
(130), slow traversal (258), and waypoint scanning (404). Testing uses a
separate fast traversal (102) and trimmed yaw scan (71).

\begin{table}[H]
\centering
\small
\caption{\textbf{Spot-with-arm window dataset.}
Four complete bags supply training windows and two separate bags supply test windows, after filtering unusable or calibration-test segments. No bag crosses the split.}
\label{tab:spot-dataset-arxiv}
\begin{tabular}{@{}l r@{}}
\toprule
\textbf{Quantity} & \textbf{Value} \\
\midrule
usable teleoperated episodes & 6 \\
typed possible-next-behavior windows & 1,070 \\
training windows & 897 \\
held-out windows & 173 \\
training / held-out bags & 4 / 2 \\
history length & 2 s \\
future-consequence horizon & 5 s \\
window stride & 1 s \\
reported query inputs & history + numeric behavior descriptors \\
\bottomrule
\end{tabular}
\end{table}

Table~\ref{tab:experiment-scope-arxiv} distinguishes this real-log study from
the cross-robot and recovery studies in
Appendices~\ref{app:accuracy_transfer} and~\ref{app:robot-decision-interface}.
The Spot row is real-robot evidence for telemetry-to-action-card prediction under logged replay; the Isaac transfer and push rows carry the branchable PHASTCore rollout evidence.

\begin{table}[H]
\centering
\scriptsize
\caption{\textbf{Benchmark roles and hierarchy for the shared query interface.}
Each artifact stresses a different part of the C-PHAST claim: simulated branch execution, typed physical-critic scoring, or real robot telemetry-to-action-card forecasting.}
\label{tab:experiment-scope-arxiv}
\setlength{\tabcolsep}{3pt}
\begin{tabularx}{\linewidth}{@{}>{\raggedright\arraybackslash}p{0.14\linewidth}
>{\raggedright\arraybackslash}p{0.18\linewidth}
>{\raggedright\arraybackslash}p{0.34\linewidth}
>{\centering\arraybackslash}p{0.08\linewidth}
>{\raggedright\arraybackslash}X@{}}
\toprule
\textbf{Experiment} & \textbf{Query mode} & \textbf{Hierarchy used} & \textbf{Real robot} & \textbf{Validated role} \\
\midrule
Isaac transfer & branchable q-only rollout & four 3-DOF leg PHASTCore rollout blocks (HAA/HFE/KFE); floating base excluded & no & cross-robot rollout accuracy and pH/compositional transfer \\
Push critic & full-state rollout plus critic & separately trained body-plus-legs C-PHAST runtime; root/joint telemetry, contact/slip geometry, and terrain context condition candidate forecasts & no & typed branch scoring, early warning, and single-decision command selection in simulation \\
Spot shadow mode & logged replay windows & interface blocks for base, leg support/contact, arm, camera/depth, and battery/power; target hierarchy for follow-up closed-loop hardware rollout & yes & real telemetry-to-action-card forecasting \\
\bottomrule
\end{tabularx}
\end{table}

\paragraph{Consequence predictors.}
\label{app:spot-models}
The Spot hierarchy organizes telemetry and outputs, not compositional dynamics.
Two direct predictors instantiate this interface. For objective-score evaluation,
multioutput ridge regression maps 28 history summaries and ten numeric behavior
descriptors to 38 consequence channels. Inputs and targets are standardized;
regularized least squares uses strength 1.0 and an unpenalized intercept.
Predicted channels are then aggregated into objective scores, rather than
learning a different physical predictor for every objective.
Panel B of Table~\ref{tab:spot-realworld-main} uses ridge regression throughout:
its first four rows share the four-bag/two-bag split, including the comparator
that predicts objective scores directly. The final row refits the consequence
predictor on five bags and evaluates the sixth, recomputing standardization
and factor scales within each training fold. Correlations are averaged equally
over the six bags and five objectives.

The energetic-channel evaluation uses a separate neural ensemble. Two-layer
history and descriptor encoders use SiLU activations and layer normalization;
a joint multilayer perceptron receives both embeddings and their elementwise
product. Sigmoid outputs enforce bounded channels, while softplus enforces
nonnegative channels. Training minimizes weighted mean-squared error on scaled
consequence targets, with seven channels weighted by 1.5; predictions are
averaged over seeds 42--46. The archived launch and retained producer defaults
specify width 32, AdamW with
learning rate 0.002 and weight decay 0.001, batch size 128, and at most 700
epochs. Early stopping uses 20\% of training-partition windows and patience 80;
these overlapping validation windows are not independent episodes. Model
checkpoints, resolved run settings, and stopping epochs were not retained
with the exported ensemble predictions, so exact retraining is unverified.
Neither predictor uses a PHASTCore transition or transferred dynamics weights.

The main subtlety is the Spot input contract.
C-PHAST is not allowed to solve the task by memorizing primitive names: the reported model receives recent telemetry plus numeric behavior descriptors, not a one-hot primitive label.
Table~\ref{tab:spot-input-contract-arxiv} separates four quantities that are easy to conflate: history observed before the query, behavior descriptors that a planner could provide, replay-only descriptor fields measured from the logged future, and consequence targets that must be predicted.

\begin{table}[H]
\centering
\scriptsize
\caption{\textbf{Spot input-contract audit.}
The held-out-bag result tests the replay version of the planner-facing interface. A deployed confirm-to-execute planner would supply behavior descriptors from intended primitive parameters rather than measured future telemetry.}
\label{tab:spot-input-contract-arxiv}
\setlength{\tabcolsep}{3pt}
\begin{tabularx}{\linewidth}{@{}lXX@{}}
\toprule
\textbf{Quantity class} & \textbf{Examples in the Spot artifact} & \textbf{Use in this paper} \\
\midrule
Decision-time history & recent speed/contact, hand motion, gripper state, image/depth quality, TF/depth novelty summaries & observed before the query; valid model input \\
Command-like behavior descriptors & commanded base speed and turn rate; gripper opening; primitive prototype descriptor & intended planner inputs; logged values are used as replay proxies \\
Future-derived behavior descriptors & realized yaw change, realized hand speed, realized hand turn over the five-second future window & used only in logged replay; must be supplied by a primitive library or forecasted before live use \\
Predicted consequence targets & work, support/contact, battery draw, collision/contact proxy, sensing quality, depth novelty, residual-like terms & unavailable before execution; predicted by the typed consequence model \\
Visualization / metadata & behavior labels, episode type, trim markers & used for analysis and plots; not part of the main reported model input \\
\bottomrule
\end{tabularx}
\end{table}

This contract gives two related query types.
The held-out objective-score task is logged-window consequence prediction: each dataset row uses the executed teleoperation future as supervision.
The same-history action-card query is a model query: it fixes one observed history and swaps in four prototype behavior descriptors, namely fast base loop, slow base loop, move-and-scan, and static arm scan, estimated from the training episodes.
That query demonstrates objective switching from one state-history context; Isaac experiments provide matched-state branch execution when actual alternatives must be run.

\paragraph{Consequences and objective scores.}
For the objective-score experiment, the regression predictor estimates
consequence channels from history summaries and behavior descriptors.
Table~\ref{tab:spot-factors-arxiv} gives their definitions and decision roles.
The predicted channels are aggregated and scaled to form a factor vector,
reused for five objectives: efficient motion, safe motion, mapping, balanced
operation, and inspection. This is the separation illustrated in
Figure~\ref{fig:signals-to-action-cards}, with a direct consequence predictor
rather than a latent dynamics rollout.

\begin{table}[H]
\centering
\scriptsize
\caption{\textbf{Spot observable-to-factor examples.}
These definitions specify telemetry-derived targets and operational proxies
for the direct consequence predictor; the objective scores use its predictions.
Objective scores scale each planner-facing factor by its training-split 10--90 percentile range before applying objective weights.}
\label{tab:spot-factors-arxiv}
\begin{tabularx}{\linewidth}{@{}lXl@{}}
\toprule
\textbf{Factor family} & \textbf{Consequence definition} & \textbf{Decision role} \\
\midrule
Progress & future base path length in the horizontal plane & reward for traversal \\
Joint work & $\sum_t\sum_j\max(\tau_{j,t}\dot q_{j,t},0)\Delta t$; absolute work replaces each positive part by $|\tau_{j,t}\dot q_{j,t}|$; both in joules & physical cost \\
Battery draw & $\max(-\overline{I}_{\mathrm{battery}},0)$, where the bar averages the summed battery current over the future horizon, in amperes & discharge-current cost \\
Support/contact & support-loss fraction, full-support fraction, contact-count deficit, foot/contact proxies & safety envelope and warning \\
Wrench/load & mean force and torque norms from available wrench or effort channels & contact/load cost \\
Collision/contact residual & hand or contact-envelope violation fraction; residual-like contact mismatch & veto or warning \\
Motion/aggression & speed, yaw motion, commanded speed and turn magnitude & blur risk, support risk, and progress estimate \\
Sensing quality & image frame/camera counts, sharpness, contrast, brightness, entropy, edge density & inspection quality \\
Depth/TF novelty & valid depth fraction, point count, observed/new voxel counts, novelty fraction, coverage rate & map gain \\
\bottomrule
\end{tabularx}
\end{table}

The reported Spot objectives apply soft weights to these factors
(Table~\ref{tab:spot-objective-weights-arxiv}). A hard veto would require an
explicit admission rule, rather than merely a larger penalty weight.

The five objectives used in the Spot result are fixed linear consumers of the planner-factor vector.
Efficient motion rewards progress while penalizing work, battery draw, command aggression, support loss, and collision.
Safe motion places stronger penalties on support loss, collision, force, and aggressive commands.
Mapping rewards progress, yaw/hand motion, image quality, and depth novelty while penalizing speed, work, battery draw, support loss, and collision.
Balanced operation trades progress, sensing motion, energy, support, force, and collision.
Inspection emphasizes image and depth quality, stable viewpoint changes, hand motion, and hand turn, with penalties for speed, near-depth risk, support loss, collision, and work.
Changing objectives therefore changes the weights applied to the same typed factor vector rather than retraining the dynamics or consequence model.

\begin{table}[H]
\centering
\scriptsize
\caption{\textbf{Objective weights for Spot action cards.}
After quantile scaling, each objective is a transparent weighting of typed planner factors.
Positive entries reward a predicted consequence, while negative entries penalize cost or risk.}
\label{tab:spot-objective-weights-arxiv}
\setlength{\tabcolsep}{3pt}
\renewcommand{\arraystretch}{1.12}
\begin{tabularx}{\linewidth}{@{}>{\raggedright\arraybackslash}p{0.13\linewidth}>{\raggedright\arraybackslash}p{0.24\linewidth}>{\raggedright\arraybackslash}p{0.34\linewidth}>{\raggedright\arraybackslash}X@{}}
\toprule
\textbf{Objective} & \textbf{Rewards} & \textbf{Penalties / veto pressure} & \textbf{Intended action card} \\
\midrule
Efficient &
progress $+1.0$ &
support loss $-1.0$, collision $-1.0$, command aggression $-0.8$, work $-0.5$, energy/progress $-0.4$, battery draw $-0.3$ &
move if progress is worth the physical cost \\
Safe &
progress $+0.2$ &
support loss $-2.0$, collision $-1.5$, force $-0.6$, command aggression $-0.5$ &
stay inside the support/contact envelope \\
Mapping &
progress $+0.7$; yaw/hand motion $+0.6/+0.8$; image/depth novelty and quality $+0.3$--$+0.6$ &
support loss $-1.0$, collision $-0.8$, speed $-0.5$, work/battery $-0.3$ &
prefer informative viewpoints without unsafe motion \\
Balanced &
progress $+0.7$; yaw/hand motion $+0.3/+0.4$ &
support loss $-1.2$, collision $-0.8$, force $-0.6$, work $-0.4$, energy/progress $-0.25$ &
trade progress, sensing, and physical risk \\
Inspection &
image/depth quality $+0.4$--$+0.9$; hand motion $+1.0$; yaw motion $+0.8$; progress $+0.25$ &
support loss $-1.2$, collision $-1.0$, speed $-0.8$, near-depth risk $-0.2$, work $-0.15$ &
slow local scan with high visual/depth fidelity \\
\bottomrule
\end{tabularx}
\end{table}

Table~\ref{tab:spot-realworld-main} gives the compact real-robot evidence ledger in the main text.
The appendix expands that ledger in five steps: Table~\ref{tab:spot-energetic-arxiv} reports direct energetic-target agreement, Table~\ref{tab:spot-objectives-arxiv} reports the held-out objective-score split, Table~\ref{tab:spot-lobo-arxiv} tests leave-one-episode-out generalization, Table~\ref{tab:spot-contract-arxiv} checks which query inputs carry signal, and Table~\ref{tab:spot-typed-direct-arxiv} compares typed readouts against direct scalar-objective training.

\begin{table}[H]
\centering
\small
\caption{\textbf{Direct energetic-target agreement on 173 held-out Spot windows.}
Targets are computed from logged joint torque/velocity or battery telemetry and predicted by the five-seed typed query-model probe.
They are telemetry-grounded physical quantities, not additional terms in the latent pH power identity.}
\label{tab:spot-energetic-arxiv}
\begin{tabular}{lcc}
\toprule
\textbf{Target} & \textbf{Pearson} & $\mathbf{R^2}$ \\
\midrule
Positive joint work & 0.973 & 0.944 \\
Absolute joint work & 0.975 & 0.947 \\
Locomotion positive work & 0.971 & 0.939 \\
Battery draw & 0.943 & 0.878 \\
\bottomrule
\end{tabular}
\end{table}

Positive-work Pearson correlation is $0.975$ on the 102 traversal windows and
$0.069$ on the 71 yaw-scan windows. Their RMSEs are similar (28.86 and 27.55\,J),
but the yaw-scan targets vary much less; pooled agreement therefore does not
establish strong within-episode tracking in both behaviors.

\begin{table}[H]
\centering
\scriptsize
\caption{\textbf{Held-out objective-score agreement on Spot logs.}
All rows use recent history plus numeric behavior descriptors without primitive identity labels.
Brackets give 2.5th--97.5th percentiles from 5,000 bootstrap resamples of
individual prediction--target pairs, with the model and normalization fixed.
Because windows overlap, these intervals describe variation within the two
test episodes, not 95\% coverage for new-episode performance.
Highest correlations are bold green; MAE scales differ across objectives.}
\label{tab:spot-objectives-arxiv}
\begin{tabular}{lccc}
\toprule
\textbf{Objective} & \textbf{Pearson} & \textbf{Spearman} & \textbf{MAE} \\
\midrule
Efficient & $0.744$ [$0.682,0.803$] & $0.753$ [$0.681,0.811$] & 0.213 \\
Safe & $0.907$ [$0.873,0.935$] & $0.875$ [$0.833,0.903$] & 0.106 \\
Mapping & $0.919$ [$0.894,0.939$] & $0.900$ [$0.861,0.926$] & 0.370 \\
Balanced & \goodcell{\textbf{$0.948$ [$0.931,0.962$]}} & \goodcell{\textbf{$0.939$ [$0.901,0.961$]}} & 0.140 \\
Inspection & $0.884$ [$0.846,0.914$] & $0.836$ [$0.767,0.886$] & 0.540 \\
\bottomrule
\end{tabular}
\end{table}

\begin{table}[H]
\centering
\small
\caption{\textbf{Leave-one-episode-out Spot stress test.}
Each of six folds trains on five episodes and holds out the sixth; metrics
are then averaged equally across episodes. Unlike the pooled fixed-split
result, every episode is evaluated out of sample. Both protocols hold out
complete episodes, but differ in coverage and aggregation.
Highest correlations are bold green; MAE scales differ across objectives.}
\label{tab:spot-lobo-arxiv}
\begin{tabular}{lccc}
\toprule
\textbf{Objective} & \textbf{Pearson} & \textbf{Spearman} & \textbf{MAE} \\
\midrule
Efficient & 0.538 & 0.412 & 0.324 \\
Safe & 0.641 & 0.590 & 0.382 \\
Mapping & \goodcell{\textbf{0.750}} & \goodcell{\textbf{0.617}} & 0.318 \\
Balanced & 0.644 & 0.517 & 0.317 \\
Inspection & 0.686 & 0.590 & 0.473 \\
\midrule
Mean & 0.652 & 0.545 & 0.363 \\
\bottomrule
\end{tabular}
\end{table}

\begin{table}[H]
\centering
\small
\caption{\textbf{Query-input ablation.}
Both current history and behavior descriptors carry signal; primitive identity alone does not.
Best per metric is shown in \textbf{bold green}.}
\label{tab:spot-contract-arxiv}
\begin{tabular}{lccc}
\toprule
\textbf{Inputs} & \textbf{Mean Pearson} & \textbf{Mean Spearman} & \textbf{Mean MAE} \\
\midrule
history only & 0.767 & 0.737 & 0.360 \\
behavior descriptors only & 0.780 & 0.778 & 0.334 \\
history + behavior descriptors & \goodcell{\textbf{0.880}} & \goodcell{\textbf{0.860}} & \goodcell{\textbf{0.274}} \\
primitive label only & -0.120 & -0.054 & 0.740 \\
history + descriptors + label & 0.854 & 0.844 & 0.313 \\
\bottomrule
\end{tabular}
\end{table}

\begin{table}[H]
\centering
\small
\caption{\textbf{Typed readouts versus direct scalar objectives.}
Typed consequences preserve scalar ranking quality while exposing the reason for the ranking.
Best per metric is shown in \textbf{bold green}.}
\label{tab:spot-typed-direct-arxiv}
\begin{tabular}{lccc}
\toprule
\textbf{Training target} & \textbf{Mean Pearson} & \textbf{Mean Spearman} & \textbf{Mean MAE} \\
\midrule
typed consequences then score & \goodcell{\textbf{0.880}} & 0.860 & \goodcell{\textbf{0.274}} \\
direct scalar objective & 0.871 & \goodcell{\textbf{0.875}} & 0.330 \\
\bottomrule
\end{tabular}
\end{table}

For a fixed observed history, the ranked future changes when the objective changes.
In the representative query in Figure~\ref{fig:spot-real-query-context}, efficient and safe objectives choose a slow-loop candidate, while mapping, balanced, and inspection objectives choose a move-and-scan candidate.
This is the operational reason to expose typed factors rather than a single learned score: a planner or operator can change the task objective while reusing the same consequence predictions.
